 \documentclass{exam}
\usepackage[spanish,activeacute]{babel}
\usepackage[utf8]{inputenc}
\usepackage[T1]{fontenc}
\usepackage[newcommands]{ragged2e}

\usepackage{
    amsmath,
    amssymb,
    eso-pic,
    float,
    graphicx,
    lmodern,
    wrapfig,
    tabularx,
    multicol,
    multirow,
    color,
    colortbl,
    lastpage,
    titlesec,
    sectsty
}

\definecolor{azul}{RGB}{33,127,190}
\sectionfont{\color{azul}}
\subsectionfont{\color{azul}}
\renewcommand{\familydefault}{\sfdefault}

\footer{}{\thepage}{}

\makeatother

\title{\LARGE\color{azul}\textbf{ICI 202 Programaci\'on 2 - Ejercicio sumativo 01 }}
\author{\normalsize \color{gray}{Prof.} \color{black}{\textbf{Ismael Figueroa, Eduardo Godoy}}}
\date{\normalsize \em \today}

\begin{document}

\AddToShipoutPictureBG*{%
  \AtPageUpperLeft{\raisebox{-\height}{\includegraphics[scale=.95]{base/header.png}}}}

\maketitle

\begin{multicols}{2} \begin{flushleft} \textbf{Nombre:} \\ \vspace*{2mm} \textbf{Rut:} \\ \vspace*{2mm} \textbf{Paralelo:} \end{flushleft} \begin{center} \begin{table}[H] \begin{tabular}{p{4cm}|p{3cm}|} \arrayrulecolor{gray!50}\cline{2-2} ~ & {\em {\scriptsize \color{gray!50}{Puntaje:}}} \\ & ~ \\ ~ & \textbf{Nota:} \\ & ~ \\ \arrayrulecolor{gray!50}\cline{2-2} \end{tabular} \end{table} \end{center} \end{multicols}

\vspace*{-18mm}
\noindent
\textbf{\\Instrucciones:}
\begin{itemize}
    \item[-] El puntaje m\'aximo  es 100 puntos.
    \item[-] Tiempo m\'aximo: 90 minutos.
    \item[-] El trabajo es \underline{\textbf{individual}}. Cualquier intento de copia, ser\'a sancionado seg\'un dicta el reglamento de la carrera.
    \item[-] Se debe subir al aula virtual un archivo comprimido con el siguiente formato \\ {Ej01\_<NombreApellidoEstudiante\_rut\_paralelo>.zip, dentro deben estar los c\'odigos fuentes requeridos.}
\end{itemize}

\noindent
\textbf{Resultados de aprendizaje a evaluar:}
\begin{enumerate}
  \item Caracter\'isticas de Orientaci\'on a objetos.
  \item Modelamiento y su implementaci\'on.
\end{enumerate}
\vspace{2mm}

\noindent
\textbf{Contenido:} Este certamen eval\'ua los siguientes temas:

\vspace{-2mm}
\begin{table}[H]
\begin{tabular}{
    !{\color{gray!50}\vrule}l
    !{\color{gray!50}\vrule}c
    !{\color{gray!50}\vrule}c
    !{\color{gray!50}\vrule}} \arrayrulecolor{gray!50} \hline
    \multicolumn{1}{!{\color{gray!50}\vrule}c}{\multirow{2}{*}{\textbf{
        Tema
    }}} &
    \multicolumn{2}{!{\color{gray!50}\vrule}c!{\color{gray!50}\vrule}}{\textbf{
        Puntajes
    }} \\ \arrayrulecolor{gray!50}\cline{2-3} &
    \multicolumn{1}{!{\color{gray!50}\vrule}c!{\color{gray!50}\vrule}}{\textbf{
        Total
    }} &
    \multicolumn{1}{c!{\color{gray!50}\vrule}}{\textbf{
        Obtenido
    }} \\ \arrayrulecolor{gray!50} \hline
    Modelamiento del problema y ejemplo de una instancia de clase.
    & \multicolumn{1}{!{\color{gray!50}\vrule}c!{\color{gray!50}\vrule}}{\textbf{
        40 pts.
    }} & \\ \arrayrulecolor{gray!50} \hline
    Ejemplo de una instancia de clase.
    & \multicolumn{1}{!{\color{gray!50}\vrule}c!{\color{gray!50}\vrule}}{\textbf{
        10 pts.
    }} & \\ \arrayrulecolor{gray!50} \hline
    Codificaci\'on de Clase con atributos y m\'etodos en leguaje java.
    & \multicolumn{1}{!{\color{gray!50}\vrule}c!{\color{gray!50}\vrule}}{\textbf{
        40 pts.
    }} & \\ \arrayrulecolor{gray!50} \hline
    Cumple con el formato y compilaci\'on de cada clase de forma correcta
    & \multicolumn{1}{!{\color{gray!50}\vrule}c!{\color{gray!50}\vrule}}{\textbf{
        10 pts.
    }} & \\ \arrayrulecolor{gray!50} \hline

\end{tabular}
\end{table}

\vspace{-7mm}
\section{\textbf{Modelamiento en Orientaci\'on a Objetos (100 pts)}}
\noindent
\textbf{Plantamiento de problema: }

\begin{questions}

  \begin{enumerate}
    \item \emph{70pts.} Cree un programa que permita realizar una correcta gestion de bodega de un antiguo almacen, para esto se requiere  implementar lo siguiente.
    \begin{itemize}

    \item \emph{(10pts)} Crear La clase  \textbf{Producto} que tenga los siguientes atributos:
      \begin{enumerate}
        \item c\'odigo: de tipo entero, cuya responsabilidad es identificar al producto.
    	  \item nombre: de tipo string, que contiene al nombre del producto.
    	  \item stock: de tipo entero, dato encargado de manejar la cantidad de productos del mismo nombre que actualmente se tiene en el almacen.
    	  \item precio: precio del producto con iva incluido.
    \end{enumerate}
    Adem\'as debe contener los m\'etodos get y set asociados a cada atributo.

    \item \emph{(5pts)} Crear la Clase \textbf{GestionBodega} que posea como atributo:
     un arreglo de tipo dinámico de tipo Producto y llamado \textbf{listaProductos} que permita almacenar los productos que crear\'a el usuario.

    La clase GestionBodega debe poseer los siguientes m\'etodos
    \begin{enumerate}

      \item \emph{(15pts)} Codificar el m\'etodo \textbf{crearProducto} que permita crear un  producto específico y asignarle un stock inicial. Se debe considerar que al asignarle el precio al producto, el sistema debe agregarle  de forma automática el iva (19\%) sobre el precio ingresado. luego de esto se procederá a guardar el resultado en el atributo precio de la clase producto. Para finalizar agregandolo al arreglo listaProductos.
      La acci\'on de crear productos debe repetirse mientras el usuario lo desee.

      \item \emph{(10pts)} Codificar el m\'etodo \textbf{listarStock} que permita visualizar el c\'odigo, nombre y stock de cada producto.


    \item \emph{(5pts)} Resuelve el problema utilizando el paradigma de Orientaci\'on a Objectos.

    \end{itemize}
    \end{enumerate}
\end{questions}

\begin{table}[H]
\centering
\begin{tabular}{
!{\color{gray!50}\vrule}p{3.9cm}
!{\color{gray!50}\vrule}p{3.6cm}
!{\color{gray!50}\vrule}p{3.6cm}
!{\color{gray!50}\vrule}p{3.6cm}
!{\color{gray!50}\vrule}} \arrayrulecolor{gray!50} \hline
    \multicolumn{4}{!{\color{gray!50}\vrule}c!{\color{gray!50}\vrule}}{\textbf{?`C\'omo ser{\'e} evaluado en este trabajo?}} \\ \arrayrulecolor{gray!50} \hline
    \textbf{\'Item} & \textbf{Logrado} & \textbf{Suficiente} & \textbf{No Logrado}\\ \arrayrulecolor{gray!50} \hline
  \newline Creaci\'on de Clase con atributos. &
  \newline  Aplica de forma correcta 40\%  ... &
  \newline Aplica parcialmente con menos de 2 errores 25\% ... &
    \newline Aplica de forma incorrecta con 3 errores o m\'as 0-9\%\\ \arrayrulecolor{gray!50} \hline
    \newline  Implementaci\'on de m\'etodos (comportamientos) asociados y correcto funcionamiento. &
    \newline  Aplica de forma correcta 10\%  ... &
    \newline Aplica parcialmente con menos de 2 errores 5\%... &
    \newline Aplica de forma incorrecta con 3 errores o m\'as 0\%\\ \arrayrulecolor{gray!50} \hline
  \newline  Implementar clase instanciadora con m\'etodo main y crea instancias de objetos. &
  \newline  Aplica de forma correcta 40\%... &
    \newline  Aplica parcialmente con menos de 2 errores 25\%... &
    \newline Aplica de forma incorrecta con 3 errores o m\'as 0-9\%\\ \arrayrulecolor{gray!50} \hline
    \newline  Cumple con el formato y compilaci\'on correcta &
    \newline  Aplica de forma correcta 10\% ... &
    \newline Aplica parcialmente con menos de x errores 5\%... &
    \newline Aplica de forma incorrecta con x errores o m\'as 0-4\%\\ \arrayrulecolor{gray!50} \hline
    Total de la secci\'on &  100\% & 60\% & 0-22\%\\ \arrayrulecolor{gray!50} \hline
\end{tabular}
\label{tbl:1}
\end{table}
\vspace{-5mm}
\textbf{Nota:} En caso de que el {\'i}tem no est{\'e} presente, tiene ponderaci{\'o}n cero.

\end{document}
